%% Journal of Econometrics submission, Elsevier elsarticle (Harvard/authoryear).
%% Template: elsarticle-template-harv.tex (elsarticle 3.x).
\documentclass[preprint,12pt,authoryear,times]{elsarticle}

\usepackage{amssymb,mathtools,amsthm,booktabs,float,siunitx}
\usepackage[hidelinks]{hyperref}
\mathtoolsset{showonlyrefs}

\journal{Journal of Econometrics}

\newtheorem{theorem}{Theorem}
\newtheorem{remark}{Remark}

\newcommand{\E}{\mathbb{E}}
\newcommand{\Q}{\mathbb{Q}}
\newcommand{\R}{\mathbb{R}}

\begin{document}

\begin{frontmatter}

\title{Identifying Risk-Neutral Densities from an Empirical Measure of Listed Quotes}

\author[ifz]{Simon A.\ Broda\corref{cor1}}
\ead{simon.broda@hslu.ch}
\affiliation[ifz]{organization={Institute of Financial Services IFZ, Lucerne University of Applied Sciences},
            addressline={Suurstoffi 1},
            city={Rotkreuz},
            postcode={6343},
            country={Switzerland}}
\cortext[cor1]{Corresponding author.}

\begin{abstract}
European option prices of a single maturity encode the risk-neutral density of the underlying at expiry. Differentiating those prices twice is ill-posed on a finite strike list. This paper rescales the calls so that they become a complementary distribution function, replaces the empirical jump measure by a Gaussian kernel, and recovers the density from the associated characteristic function in closed form. The estimator is a finite Gaussian mixture, one term per strike. Completing the unquoted tails, placing masses at cell centers, and combining two kernel widths are quadrature corrections; the resulting estimator is in-fill consistent. The same kernel supplies a smooth estimate of the density and a call interpolant. On listed S\&P~500 and Nasdaq-100 chains, at a density-scale bandwidth, that interpolant reprices out-of-the-money mids more closely than the positive convolution approximation of \citet{bondarenko2003}, the constrained least squares of \citet{yatchew2006}, the shape-constrained local polynomial of \citet{aitsahalia2003}, and a kernel of the call of \citet{priestley1972}, including on an even/odd hold-out averaged over both directions. 
\end{abstract}

\begin{keyword}
risk-neutral density \sep option prices \sep kernel estimation \sep in-fill consistency \sep state-price density
\MSC[2010] 62G07 \sep 62G20 \sep 62P05 \sep 91G20
\JEL C14 \sep C58 \sep G13
\end{keyword}

\end{frontmatter}

\section{Introduction}
\label{sec:intro}

\citet{breeden1978} identified the risk-neutral density of the underlying at a fixed expiry with the second strike derivative of the corresponding European call, scaled by a discount factor. The identity assumes a twice differentiable price curve. A market supplies a finite ordered list of strikes and a quote at each knot. Numerical second differentiation of that list is ill-posed: interpolants oscillate, change sign, and depend on the interpolating scheme. The literature responds with parametric mixtures and Gram--Charlier expansions, implied-volatility interpolants, local polynomials of the call, and \citet{bondarenko2003}'s positive convolution approximation, which forces the density to be a mixture of a positive kernel with a discrete probability fitted to prices. Section~\ref{sec:lit} reviews these approaches.

This paper estimates the density from a single-maturity strip by rescaling the calls to a complementary distribution function, replacing the empirical jump measure by a Gaussian kernel, and inverting the associated characteristic function. For a Gaussian the inversion is a finite sum, one term per strike. Completing the unquoted tails, placing each mass at the center of its cell, and combining two kernel widths are quadrature corrections. As the strike mesh fills a compact interval and the bandwidth shrinks more slowly than the mesh, the estimator converges to the true density on the interior of that interval.

The usual split in this literature is that interpolants of the call or of implied volatility reprice listed quotes well and do not yield a usable density, while methods that are densities by construction cannot interpolate a dense, slightly arbitrageable chain. \citeauthor{bondarenko2003}'s method is a sieve of fixed dimension: it does not become finer as more strikes arrive, and it is a density by construction. \citet{yatchew2006} replace the quoted calls by the nearest decreasing convex list; the second differences of that list are a discrete nonnegative measure, not a curve. \citet{aitsahalia2003} smooth the same constrained list with a local polynomial, so the density is nonnegative by design and the call cannot match quotes that violate convexity. The estimator of this paper returns a closed-form smooth estimate of \(f_{\Q}\) and, from the same kernel, a call interpolant. On listed S\&P~500 and Nasdaq-100 chains, at a density-scale bandwidth and after the two-width combination, that interpolant has smaller out-of-the-money root-mean-square error (RMSE) than the no-arbitrage projection, than \citeauthor{bondarenko2003} and \citeauthor{aitsahalia2003}, and than a \citet{priestley1972} kernel of the call, in sample and on an even/odd hold-out averaged over both directions. It can lie below the \citeauthor{yatchew2006} floor because it is not required to be convex. 

The remainder of this manuscript is structured as follows. Section~\ref{sec:lit} reviews related work. Section~\ref{sec:est} derives the characteristic function of the normalized call and the closed-form estimator, and compares it to kernel smoothing of the call. Section~\ref{sec:refine} records three refinements. Section~\ref{sec:h} chooses the bandwidth. Section~\ref{sec:ifc} proves in-fill consistency and illustrates it on exact Heston and variance-gamma quotes. Section~\ref{sec:emp} presents an empirical illustration that evaluates our estimator against four alternative methods on four listed chains. Section~\ref{sec:conc} concludes.

\section{Related Work}
\label{sec:lit}

The identity of Section~\ref{sec:est} is the continuous-strike version of a spanning argument of \citet{ross1976}: if one could trade a put at every strike, the second derivative of the put curve would be a density, and one could synthesize any sufficiently regular payoff. \citet{dupire1994} obtained the analogous result in calendar time, which produces a local-volatility equation. Observed markets supply neither a continuum of strikes nor noiseless quotes.

A first family of methods specifies the density of \(S_T\) as a parametric curve and fits the implied option prices to the quotes. \citet{ritchey1990} and \citet{melick1997} take the density to be a mixture of lognormals, so that the call price is a mixture of Black--Scholes prices; \citet{bahra1997} and \citet{soderlind1997} use related specifications for interest-rate and foreign-exchange options. Two or three components typically suffice to generate a unimodal, skewed density, but the fit is poorly identified once more than two components are admitted, and the approximation to a Heston or jump-diffusion density need not vanish as more quotes arrive. Expansion methods perturb a baseline density---usually lognormal---by Gram--Charlier or Hermite terms \citep{jarrowrudd1982,corradosu1996,rompolis2007,xiu2014}. The expansions can match a finite number of moments, but they are not guaranteed to stay nonnegative, and the truncated series can oscillate in the tails.

A second, and in applied work the most common, family interpolates the Black--Scholes implied-volatility function of strike and converts the interpolant back to prices before differentiating twice. \citet{shimko1993} fitted a quadratic in strike; \citet{malz1997} used a cubic spline in delta; \citet{bliss2002} a smoothing spline. \citet{figlewski2010} completes the tails with a generalized extreme-value distribution. Because implied volatility is a much flatter function of strike than the call itself, a low-order interpolant often has small bias when quotes are accurate. The drawback is that implied volatility is a nonlinear transformation that inflates price noise where the derivative of price with respect to implied volatility is small, i.e.\ in the tails. An interpolating spline then treats those noisy tail quotes as exact, and the second derivative of the resulting call curve oscillates and can become negative. Shape-constrained smoothers of the implied-volatility surface \citep{fengler2005,fengler2009,gatheraljacquier2014,ludwig2015} address the sign problem at the cost of a more involved optimization.

A third family estimates the call as a function of strike by kernel regression or local polynomials and obtains the second derivative from the fit. \citet{priestley1972} fitted a kernel to observed function values at the design points; the risk-neutral density is obtained by differentiating that fit twice. \citet{aitsahalia1998} estimated the call as a function of \((S,K,T,r,q)\) by kernel regression; \citet{aitsahalia2000} used the estimator to recover an implied risk-aversion function. The method is data intensive and, in the original application, pooled option prices across many trading dates. \citet{aitsahalia2003} developed a local polynomial estimator that imposes monotonicity and convexity of \(C\) in \(K\), hence nonnegativity of the density. Related constrained least-squares estimators are studied by \citet{yatchew2006} and \citet{birke2009}; \citet{grith2012} survey the kernel literature. \citet{fanmancini2009} combine a parametric guide with a nonparametric correction. In each case the density is the second strike derivative of a smoother of \(C\). The estimator of Section~\ref{sec:est} rescales the calls to a complementary cumulative distribution function (cdf) and applies the kernel to the empirical jump measure, so the formula is a first derivative of a Gaussian kernel density and is a finite sum. Section~\ref{sec:naive} records the comparison. Section~\ref{sec:emp} implements the cubic \citeauthor{priestley1972} estimator of that section and the two-step procedure of \citeauthor{aitsahalia2003}: a least-squares projection of the quoted calls onto the decreasing convex cone, followed by a local polynomial of the constrained values.

\citet{bondarenko2003} proposed the positive convolution approximation (PCA). The density is required to lie in the set of functions obtained by convolving a discrete probability with a positive kernel, and the discrete probability is chosen to minimize the sum of squared pricing errors. The kernel enforces smoothness and, together with positive weights, guarantees that the result is a density. In \citeauthor{bondarenko2003}'s comparison, the method dominated mixtures of lognormals, Hermite expansions, and several regularization methods. Subsequent simulation evidence on affine jump-diffusions is consistent with this ranking \citep{lu2019}. Section~\ref{sec:emp} implements a log-price analog: the kernel is Gaussian on \(\log S\), so that the density is a mixture of lognormals sharing a common width, fitted by nonnegative least squares. The same section implements the interpolating case of \citeauthor{yatchew2006}.

\citet{rubinstein1994} and \citet{jackwerth1996} recovered probabilities on a binomial tree by minimizing a distance to a prior subject to pricing constraints. \citet{stutzer1996} and \citet{buchen1996} maximized entropy. These methods produce a valid density by construction, but the finite-support discretization and the choice of prior affect the tails. Viewing the problem as a linear inverse problem, \citet{cont2004} calibrated jump-diffusion models under a relative-entropy penalty; \citet{monnier2014} used the singular value decomposition of a restricted put operator.

A number of papers compare methods on market data, where the true density is unknown, so that the criterion is typically pricing error or stability over time \citep{cooper1999,jondeau2000,coutant2001,anagnou2005,bliss2002,liu2007}. Studies that generate options from a known process are less common. \citet{bondarenko2003} finds that his method outperforms mixtures, expansions, and several regularization methods on integrated squared error. \citet{lai2014} finds that kernel regression outperforms spline interpolation and neural nets on pricing error. \citet{santos2015} compare mixtures, Edgeworth expansions, hypergeometric functions, and a \citeauthor{bliss2002} smoothing spline against Heston and other targets. \citet{lu2019} compares methods on affine jump-diffusions and again finds \citeauthor{bondarenko2003}'s method accurate. The typical ranking is that implied-volatility interpolants price well when quotes are clean, that mixtures are unstable once more than two components are admitted, and that \citeauthor{bondarenko2003}'s method is the most accurate among methods that do not use a parametric density of \(S_T\), when the true density is known. Section~\ref{sec:ifc} is in this vein: the targets are Heston and variance-gamma, and the true density is known. Section~\ref{sec:emp} scores listed index chains, where the criterion is pricing error, against \citeauthor{bondarenko2003}, \citeauthor{yatchew2006}, \citeauthor{aitsahalia2003}, and \citeauthor{priestley1972}.

\section{Construction of the Estimator}
\label{sec:est}

\subsection{Construction}
\label{sec:construct}

Assume a risk-neutral measure \(\Q\) under which \(S_T\) has a density \(f_{\Q}\) on \((0,\infty)\) and \(\E^{\Q}[S_T]<\infty\). Let \(r\) be the continuously compounded interest rate and \(q\) the continuous dividend yield. Write \(F=S_0e^{(r-q)T}\) for the forward. Under \(\Q\),
\begin{align}
S_0e^{-qT}
&\;=\;
e^{-rT}\,\E^{\Q}[S_T]
\nonumber\\
&\;=\;
e^{-rT}\int_{0}^{\infty}x\,f_{\Q}(x)\,\mathrm{d}x,
\label{eq:S0}
\end{align}
hence \(\E^{\Q}[S_T]=F\). Write \(C\) for the corresponding call curve:
\begin{align}
C(K)
&\;=\;
e^{-rT}\,\E^{\Q}\bigl[(S_T-K)^+\bigr]
\nonumber\\
&\;=\;
e^{-rT}\int_K^{\infty}(x-K)\,f_{\Q}(x)\,\mathrm{d}x.
\label{eq:C}
\end{align}
Now define the normalized option price
\begin{equation}
\label{eq:Ctil}
\tilde C(K)
\;\equiv\;
\frac{C(K)}{S_0e^{-qT}}
\;=\;
\frac{\displaystyle\int_K^{\infty}(x-K)\,f_{\Q}(x)\,\mathrm{d}x}{F}.
\end{equation}

The normalized call price is a decreasing function of \(K\): it is equal to \(1\) if \(K=0\) and equal to \(0\) if \(K=\infty\). Thus the normalized call price, as a function of \(K\), is a complementary cdf, say \(\tilde C(K)=1-P(K)\). The characteristic function (hereafter c.f.) of the law of \(P\) is the Fourier--Stieltjes transform
\begin{equation}
\label{eq:phi-P-stieltjes}
\varphi_P(s)
\;=\;
\int_{(0,\infty)}e^{\mathrm{i}s K}\,\mathrm{d}P(K).
\end{equation}
For a cdf on \([0,\infty)\) and \(s\neq 0\), \eqref{eq:phi-P-stieltjes} is
\begin{equation}
\varphi_P(s)
\;=\;
1+\mathrm{i}s\int_0^{\infty}e^{\mathrm{i}s K}\,\tilde C(K)\,\mathrm{d}K.
\end{equation}
Substitute \eqref{eq:Ctil}. Fubini yields
\[
\int_0^{\infty}e^{\mathrm{i}s K}\,\tilde C(K)\,\mathrm{d}K
\;=\;
\frac{1}{F}\int_0^{\infty}f_{\Q}(x)\Biggl[\int_0^x e^{\mathrm{i}s K}(x-K)\,\mathrm{d}K\Biggr]\mathrm{d}x.
\]
The inner integral is \(-x/(\mathrm{i}s)+(e^{\mathrm{i}s x}-1)/(\mathrm{i}s)^2\). Hence, for \(s\neq 0\),
\begin{align}
\varphi_P(s)
&\;=\;
1+\frac{\mathrm{i}s}{F}\int_0^{\infty}f_{\Q}(x)\Biggl[-\frac{x}{\mathrm{i}s}+\frac{e^{\mathrm{i}s x}-1}{(\mathrm{i}s)^2}\Biggr]\mathrm{d}x
\nonumber\\
&\;=\;
1-F^{-1}\E^{\Q}[S_T]+\frac{\varphi_{\Q}(s)-1}{\mathrm{i}s\,F}
\nonumber\\
&\;=\;
\frac{\varphi_{\Q}(s)-1}{\mathrm{i}s\,F},
\label{eq:phi-P}
\end{align}
where \(\varphi_{\Q}(s)=\E^{\Q}[e^{\mathrm{i}s S_T}]\) and \(\E^{\Q}[S_T]=F\) by \eqref{eq:S0}. Solving for \(\varphi_{\Q}\), one obtains
\begin{equation}
\label{eq:phi-Q}
\varphi_{\Q}(s)
\;=\;
1+\mathrm{i}s\,F\,\varphi_P(s).
\end{equation}
Read from the right, \eqref{eq:phi-P} is a pricing formula: a model that supplies \(\varphi_{\Q}\) yields \(\varphi_P\), hence the complementary cdf \(\tilde C=1-P\) and the call \citep{lewis2000}. This paper inverts the identity: \(\varphi_P\) is constructed from the quotes, and \(\varphi_{\Q}\) is recovered.

The c.f.\ \(\varphi_{\Q}\) is that of \(S_T\). Fourier inversion yields
\begin{equation}
\label{eq:f-Q}
f_{\Q}(K)
\;=\;
\frac{1}{2\pi}\int_{-\infty}^{\infty}e^{-\mathrm{i}s K}\,\varphi_{\Q}(s)\,\mathrm{d}s.
\end{equation}
Substitute \eqref{eq:phi-Q} into \eqref{eq:f-Q}:
\begin{equation}
\label{eq:f-Q-plug}
f_{\Q}(K)
\;=\;
\frac{1}{2\pi}\int_{-\infty}^{\infty}e^{-\mathrm{i}s K}\,\mathrm{d}s
+\frac{F}{2\pi}\int_{-\infty}^{\infty}e^{-\mathrm{i}s K}\,(\mathrm{i}s)\,\varphi_P(s)\,\mathrm{d}s.
\end{equation}
The first integral is a Dirac mass at \(K=0\). For \(K>0\),
\begin{equation}
\label{eq:f-Q-pos}
f_{\Q}(K)
\;=\;
\frac{F}{2\pi}\int_{-\infty}^{\infty}e^{-\mathrm{i}s K}\,(\mathrm{i}s)\,\varphi_P(s)\,\mathrm{d}s.
\end{equation}

Observed calls at strikes \(0<K_1<\cdots<K_m\) supply an empirical complementary cdf \(\tilde C_i=C_i/(S_0e^{-qT})\) and hence an empirical cdf \(P_i=1-\tilde C_i\). Write
\begin{equation}
\label{eq:nu}
\nu_m
\;=\;
\sum_{i=1}^m\Delta P_i\,\delta_{K_i},
\qquad
\Delta P_i
\;=\;
P_i-P_{i-1},
\quad
P_0:=0.
\end{equation}
The c\`adl\`ag interpolant of these knots is a step function \(P_m\). Its Fourier--Stieltjes transform is the finite sum \(\int e^{\mathrm{i}s K}\,\mathrm{d}P_m(K)=\sum_i\Delta P_i e^{\mathrm{i}s K_i}\), which is also the c.f.\ of \(\nu_m\). Replace \(\varphi_P\) by the product of that sum with the Fourier transform of a Gaussian kernel of width \(h\),
\begin{equation}
\label{eq:phi-hat}
\hat\varphi_P(s)
\;=\;
\Biggl(\sum_{i=1}^m\Delta P_i e^{\mathrm{i}s K_i}\Biggr)e^{-h^2 s^2/2}.
\end{equation}
Formally, the estimator is \eqref{eq:f-Q-pos} with \(\varphi_P\) replaced by \(\hat\varphi_P\):
\begin{equation}
\label{eq:fhat-Q-ft}
\hat f_{\Q}(K)
\;=\;
\frac{F}{2\pi}\int_{-\infty}^{\infty}e^{-\mathrm{i}s K}\,(\mathrm{i}s)\Biggl(\sum_{i=1}^m\Delta P_i e^{\mathrm{i}s K_i}\Biggr)e^{-h^2 s^2/2}\,\mathrm{d}s.
\end{equation}
The integrand is a finite sum. The Fourier transform of \(e^{\mathrm{i}s K_i}e^{-h^2 s^2/2}\) is the \(N(K_i,h^2)\) density
\begin{equation}
\label{eq:kappa}
\kappa_h(u)
\;=\;
\frac{1}{h\sqrt{2\pi}}\,\exp\Bigl(-\frac{u^2}{2h^2}\Bigr),
\end{equation}
and multiplication by \(\mathrm{i}s\) is differentiation up to sign: if \(\varphi(s)=\int e^{\mathrm{i}s u}\kappa_h(u-K_i)\,\mathrm{d}u\), then the transform of \(\kappa_h'(\,\cdot\,-K_i)\) is \(-\mathrm{i}s\,\varphi(s)\), so the transform of \(\mathrm{i}s\,\varphi(s)\) is \(-\kappa_h'(\,\cdot\,-K_i)\). Therefore
\begin{equation}
\label{eq:fhat-Q}
\hat f_{\Q}(K)
\;=\;
-F\sum_{i=1}^m\Delta P_i\,\kappa_h'(K-K_i)
\;=\;
F\sum_{i=1}^m\Delta P_i\,\frac{K-K_i}{h^2}\,\kappa_h(K-K_i),
\end{equation}
where \(\kappa_h'(u)=-(u/h^2)\kappa_h(u)\). Equation \eqref{eq:fhat-Q} is a closed-form sum. 
%It is also \(-F\) times the first derivative of a Gaussian kernel density of \(\nu_m\). 
The interpolant of the complementary-call cdf is the kernel cdf of the same measure: no derivative of the kernel is required. Let \(\mathcal{K}\) denote the standard Gaussian cdf. Then
\begin{equation}
\label{eq:P-hat}
\hat P(K)
\;=\;
\sum_{i=1}^m\Delta P_i\,\mathcal{K}\bigl((K-K_i)/h\bigr),
\end{equation}
and calls are \(\hat C=S_0e^{-qT}(1-\hat P)\). Differentiating \(\hat P\) twice and multiplying by \(-F\) recovers \eqref{eq:fhat-Q}.
\label{sec:closed}

\subsection{Comparison to Naive Kernel Smoothing}
\label{sec:naive}

In the continuum, \citeauthor{breeden1978} plus a kernel is unambiguous. If \(C\) is the true call curve,
\begin{equation}
\label{eq:smooth-C}
(C\ast\kappa_h)''
\;=\;
C''\ast\kappa_h
\;=\;
e^{-rT}\,(f_{\Q}\ast\kappa_h).
\end{equation}
Smoothing \(C\) and differentiating twice is a kernel smooth of \(f_{\Q}\). On a strike list the discretization matters.

A \citeauthor{priestley1972} estimator treats the observations as levels \(C_i\) and weights them by the design spacing \(\Delta K_i=K_i-K_{i-1}\) (\(K_0:=K_1-(K_2-K_1)\)):
\begin{equation}
\label{eq:pc}
\hat C^{\mathrm{PC}}(K)
\;=\;
\sum_{i=1}^m\Delta K_i\,C_i\,\kappa_h(K-K_i),
\qquad
\hat f_{\Q}^{\mathrm{PC}}(K)
\;=\;
e^{rT}\sum_{i=1}^m\Delta K_i\,C_i\,\kappa_h''(K-K_i),
\end{equation}
where \(\kappa_h''(u)=(u^2/h^2-1)h^{-2}\kappa_h(u)\). The equivalent kernel for the density is \(\kappa''\). The interpolant is a mixture of the prices.

That sum is a Riemann approximation to \eqref{eq:smooth-C}. It treats the observed levels as atoms of a measure. They are not: \(C\) is a smooth function of \(K\), observed at the knots. Convolution requires a function on the line. A linear interpolant of \(C\) has second derivative zero almost everywhere and Dirac masses at the kinks, so it cannot recover \(f_{\Q}\). The usual construction is a cubic spline \(C^{\mathrm{cs}}\) through the quotes, extended by linear wings of slope in \([-e^{-rT},0]\):
\begin{equation}
\label{eq:pc-cubic}
\hat f_{\Q}^{\mathrm{cs}}(K)
\;=\;
e^{rT}\,(C^{\mathrm{cs}}\ast\kappa_h)''(K).
\end{equation}
A cubic is the lowest-degree interpolant whose second derivative is not identically zero on each cell; the same degree appears in local-polynomial estimation of a second derivative \citep{aitsahalia2003}. Without the interpolant, \eqref{eq:pc} places a large weight on \(\kappa''\) wherever \(C\) is large, including deep in the money where \(C\) is nearly linear and \(C''\) is nearly zero.

Equation \eqref{eq:fhat-Q} does not need a spline. The object being smoothed is already a discrete measure, \(\nu_m=\sum_i\Delta P_i\,\delta_{K_i}\). Convolution with \(\kappa_h\) is the finite mixture \(\sum_i\Delta P_i\,\kappa_h(\,\cdot\,-K_i)\), and one derivative of that mixture is \eqref{eq:fhat-Q}. The c\`adl\`ag interpolant of \(P\) is a step function; smoothing a step function is kernel density estimation, not kernel regression of a curve.

On an equal grid with already monotone \(C\), \(\Delta C_i\approx C'(K_i)\delta\), so \eqref{eq:pc} and \eqref{eq:fhat-Q} are close. They are not close when \(C\) is large and locally linear, as a deep in-the-money call is. Section~\ref{sec:emp} implements the cubic interpolant \eqref{eq:pc-cubic} on listed chains.

\section{Three Refinements}
\label{sec:refine}

\subsection{Completing the Tails}
\label{sec:tails}

The empirical measure \eqref{eq:nu} puts mass \(P(K_1)\) at the first quoted strike and mass \(1-P(K_m)\) nowhere. Both ends are wrong as a cdf on \((0,\infty)\).

On \((0,K_1]\) a deep in-the-money call is intrinsic, \(C(K)=e^{-rT}(F-K)\), hence \(\tilde C(K)=1-K/F\) and \(P(K)=K/F\). That is a straight line through the origin. The estimator matches the observed \(P(K_1)\) rather than \(K_1/F\),
\begin{equation}
\label{eq:left}
P(K)
\;=\;
P(K_1)\,\frac{K}{K_1},
\qquad
0<K\le K_1.
\end{equation}
When the first quote is at intrinsic, \(P(K_1)=K_1/F\) and \eqref{eq:left} is exactly \(K/F\). The line is replaced by \(n_L=\mathrm{round}(K_1/\delta)\) equal jumps of size \(P(K_1)/n_L\) on a uniform grid of \((0,K_1]\), where \(\delta\) is the median quoted gap, and the original jump \(\Delta P_1=P(K_1)\) at \(K_1\) is removed. That choice continues the quoted mesh through the splice at \(K_1\).

On the right, \(\tilde C(\infty)=0\) so \(P(\infty)=1\). If \(C_m>0\), a European call still has to fall to zero. The residual call is continued linearly with slope in \([-e^{-rT},0)\) until \(C=0\) at an intercept \(K_{\mathrm{end}}\). That wing is a line. The last listed gap is much coarser than \(\delta\), so the knots are a quadrature of the line rather than a continuation of the quoted mesh. Eight knots per last gap,
\begin{equation}
\label{eq:nR}
n_R
\;=\;
\mathrm{round}\bigl(8(K_{\mathrm{end}}-K_m)/(K_m-K_{m-1})\bigr),
\end{equation}
are equally spaced in \((K_m,K_{\mathrm{end}}]\) and converted to \(P=1-\tilde C\). If \(C_m=0\), \(P(K_m)=1\) already and this wing is empty.

Only the left completion changes the estimator on the Heston and variance-gamma chains below: the first quoted call is at intrinsic, and the last quoted call is already zero (dense) or of order \(10^{-5}\) (sparse).

\subsection{Midpoint Rule}
\label{sec:mid}

Placing each jump \(\Delta P_i\) at the right endpoint of its cell \((K_{i-1},K_i]\) is a rectangle rule for the Stieltjes integral against \(P\). The midpoint rule puts the same mass at the cell center. Write \(K_0:=0\) and
\begin{equation}
\label{eq:mid}
\tilde K_j
\;=\;
\tfrac12\bigl(K_{j-1}^{\ast}+K_j^{\ast}\bigr),
\end{equation}
where \(K_1^{\ast}<\cdots<K_N^{\ast}\) is the completed strike list. For a Gaussian integrand on an equally spaced grid the midpoint error is smaller than any power of \(\delta/h\). Combined with both tails,
\begin{equation}
\label{eq:headline}
\hat f_{\Q}(K)
\;=\;
F\sum_{j=1}^{N}\Delta_j\,\frac{K-\tilde K_j}{h^2}\,\kappa_h(K-\tilde K_j),
\end{equation}
with \(\kappa_h\) as in \eqref{eq:kappa} and \(\Delta_j\) the jump on \((K_{j-1}^{\ast},K_j^{\ast}]\). Equation \eqref{eq:fhat-Q} is the same sum with \(\tilde K_j\) replaced by the right endpoints \(K_j^{\ast}\). The interpolant on that support is \eqref{eq:P-hat} with each \(K_i\) replaced by \(\tilde K_j\),
\begin{equation}
\label{eq:P-hat-mid}
\hat P(K)
\;=\;
\sum_{j=1}^{N}\Delta_j\,\mathcal{K}\bigl((K-\tilde K_j)/h\bigr).
\end{equation}

\subsection{Twicing}
\label{sec:twice}

The Gaussian kernel in \eqref{eq:headline} has second moment \(h^2\), so the leading bias of \(\hat f_{\Q}\) is \(O(h^2)\). When the quotes are treated as exact there is no observation-error variance to inflate, and a higher-order kernel is admissible. Twicing \citep{schucanysommers1977} replaces the kernel transform \(e^{-h^2 s^2/2}\) in \eqref{eq:phi-hat} by
\begin{equation}
\label{eq:phi-twice}
2e^{-h^2 s^2/2}-e^{-h^2 s^2},
\end{equation}
which is the Fourier transform of \(2\kappa_h-\kappa_{h\sqrt{2}}\). The equivalent kernel has vanishing second moment, so the \(O(h^2)\) bias becomes \(O(h^4)\). In closed form,
\begin{equation}
\label{eq:twice}
\hat f_{\Q}^{\mathrm{tw}}(K)
\;=\;
2\hat f_{\Q,h}(K)-\hat f_{\Q,h\sqrt{2}}(K),
\end{equation}
and the same combination is applied to the interpolant \eqref{eq:P-hat-mid}, \(\hat P^{\mathrm{tw}}=2\hat P_h-\hat P_{h\sqrt{2}}\). The equivalent kernel takes negative values, so \(\hat f_{\Q}^{\mathrm{tw}}\) is not a density. The estimator remains a finite Gaussian mixture.

Twicing a truncated measure amplifies leakage from the missing tails. It is used only on the completed, midpoint-placed support. Equation \eqref{eq:twice}, applied to \eqref{eq:headline}, is the estimator used below.

\section{Choice of Bandwidth}
\label{sec:h}

The width \(h\) is the only tuning constant in \eqref{eq:twice}. Two rules are used below.

With exact quotes the bias is the usual \(O(h^2)\) kernel bias, or \(O(h^4)\) after twicing, and the remaining discretization error after midpoints is smaller than any power of \(\delta/h\) on an equally spaced Gaussian grid. For a Gaussian it is enough that \(h\) be a small multiple of the typical quoted gap. Write \(\delta\) for the median of \(\{K_i-K_{i-1}\}\) on the quoted strikes, not on the filled wings, and take
\begin{equation}
\label{eq:h-mesh}
h
\;=\;
\min\bigl(1.2\,\delta,\; 0.30\,(K_m-K_1)\bigr).
\end{equation}
The cap binds only on a very short chain. On an equally spaced grid every gap equals \(\delta\). This is the rule in Section~\ref{sec:heston}, where the quotes are exact.

When quotes are noisy, \eqref{eq:h-mesh} inherits the noise in the jumps. A listed chain is a noisy sample of the call curve, not an in-fill of an exact function. The object being smoothed is then closer to a density of the jumps than to a second derivative from a vanishing mesh. Replacing the mesh by the width of \(S_T\) and using the Gaussian exponent of \citet{silverman1986}, \(n^{-1/5}\), yields
\begin{equation}
\label{eq:h-den}
h
\;=\;
1.06\,F\,\hat\sigma_{\mathrm{ATM}}\sqrt{T}\,n^{-1/5},
\end{equation}
where \(\hat\sigma_{\mathrm{ATM}}\) is the at-the-money (ATM) Black--Scholes implied volatility, the median of contracts with \(|K-F|\le 0.03F\) (or the median of all finite implied volatilities if that window is empty), and \(n\) is the number of masses after the right wing of Section~\ref{sec:tails}, not the left filler and not necessarily the number of listed quotes. The scale is the scale of \(S_T\), not the standard deviation of the jumps \(\Delta_j\). The ordinary \citeauthor{silverman1986} rule computed from those jumps oversmooths, because that standard deviation is of the order of the chain width. This is the rule in Section~\ref{sec:emp}.

Theorem~\ref{thm:ifc} requires \(h\to 0\) and \(\delta/h^{3/2}\to 0\) on a general mesh, or \(\delta/h\to 0\) on an equally spaced Gaussian grid. On a sequence of grids that fill a fixed interval with mesh \(\delta\propto n^{-1}\), the density-scale rule has \(h\propto n^{-1/5}\). Then \(h\to 0\) and \(\delta/h^{3/2}=O(n^{-7/10})\to 0\), so \eqref{eq:h-den} satisfies the theorem for exact quotes. The listed chains below are noisy, and the theorem is silent on observation error. The mesh rule keeps \(h/\delta=1.2\), so \(\delta/h\) does not tend to zero and the polynomial midpoint bound \(O(\delta^2/h^3)\) does not vanish. Equation \eqref{eq:h-mesh} is not covered by Theorem~\ref{thm:ifc}. It is a fixed-ratio quadrature for exact, equally spaced quotes, used on the chains of Section~\ref{sec:heston}.

%An even/odd split that interpolates the ecdf \(P\) on even strikes and scores the odd \(P_i\) targets \(P\), not \(P''\), and undersmooths. It is not used as a bandwidth.

\section{In-Fill Consistency}
\label{sec:ifc}

Write \(P\) for the complementary-call cdf of Section~\ref{sec:est}, so that \(f_{\Q}=-F\,P''\) on \((0,\infty)\). The midpoint estimator \eqref{eq:headline} is
\[
\hat f_{\Q}(K)
\;=\;
-F\,(\nu_N\ast\kappa_h)'(K),
\]
where \(\nu_N=\sum_j\Delta_j\delta_{\tilde K_j}\) is the completed, midpoint-placed empirical measure of \(P\).

\begin{theorem}
\label{thm:ifc}
Let \(f_{\Q}\) be four times continuously differentiable on \((0,\infty)\) with \(\E^{\Q}[S_T]<\infty\), and let \(C\) be the corresponding call curve. Fix a compact interval \([A,B]\subset(0,\infty)\). Consider a sequence of exact quote grids \(0<K_1^{(m)}<\cdots<K_m^{(m)}\) with mesh \(\delta_m=\max_i(K_i^{(m)}-K_{i-1}^{(m)})\to 0\), \(K_1^{(m)}\le A\) and \(K_m^{(m)}\ge B\). Complete the tails as in Section~\ref{sec:tails} and place masses at midpoints as in Section~\ref{sec:mid}. Let \(h_m>0\). If \(h_m\to 0\) and \(\delta_m/h_m^{3/2}\to 0\), then both the midpoint estimator \eqref{eq:headline} and its twiced version \eqref{eq:twice} converge to \(f_{\Q}\) uniformly on every compact subset of \((A,B)\). If the grid is equally spaced, the Gaussian Euler--Maclaurin argument of the proof yields the same conclusion under the weaker condition \(\delta_m/h_m\to 0\).
\end{theorem}

\begin{proof}
The midpoint estimator is \(\hat f_{\Q}=-F(\nu_N\ast\kappa_h)'\), and \(f_{\Q}=-F P''\). Split
\[
(\nu_N\ast\kappa_h)'-P''
\;=\;
\bigl((\nu_N-P)\ast\kappa_h'\bigr)
+(P'\ast\kappa_h-P')'.
\]
The second term is ordinary kernel bias of \(P''=-f_{\Q}/F\): \((P'\ast\kappa_h-P')'=(P''\ast\kappa_h)-P''\). Because \(\kappa\) is even, \(\int\kappa=1\) and \(\int u\kappa(u)\,\mathrm{d}u=0\). If \(f_{\Q}''\) is bounded on a neighborhood of \([A,B]\), the difference is \(O(h^2)\) uniformly on compact subsets of \((A,B)\) several bandwidths away from the endpoints. The twiced kernel \(2\kappa_h-\kappa_{h\sqrt{2}}\) has Fourier transform \eqref{eq:phi-twice}, hence vanishing second moment, so the same expansion starts at order \(h^4\) whenever \(f_{\Q}\) is \(C^4\).

The first term is discretization. On each completed cell the midpoint rule integrates \(\kappa_h'(K-\,\cdot\,)\) against \(\mathrm{d}P\). For a \(C^2\) integrand the error on a cell of length at most \(\delta\) is \(O(\delta^3)\) times the second derivative of the integrand. That second derivative is \(O(h^{-4})\) for \(\kappa_h'\), and \(O(h^{-4})\) as well for the twiced equivalent kernel. The kernel window contains \(O(h/\delta)\) cells, so a polynomial bound on the total discretization error is \(O(\delta^2/h^3)\), which vanishes if \(\delta/h^{3/2}\to 0\). For a Gaussian integrand on an equally spaced grid the Euler--Maclaurin formula (the Fourier transform of a Gaussian is a Gaussian) makes the midpoint error smaller than any power of \(\delta/h\), so the weaker condition \(\delta/h\to 0\) already suffices. The c\`adl\`ag placement of Section~\ref{sec:closed}, with masses at the right endpoints, is a rectangle rule: the same calculation yields an error \(O(\delta/h^3)\), which vanishes if \(\delta/h^3\to 0\).

The filled wings are supported on \((0,K_1]\) and, if \(C_m>0\), on \((K_m,K_{\mathrm{end}}]\). Their contribution at an evaluation point \(K\in[A+\varepsilon,B-\varepsilon]\) is at most a constant times \(h^{-2}\exp(-\varepsilon^2/2h^2)\), because \(K_1\le A\) and \(K_m\ge B\). That bound tends to zero faster than any power of \(h\).

Multiplying by \(-F\) yields the claim for \(\hat f_{\Q}\) and for \(\hat f_{\Q}^{\mathrm{tw}}\).
\end{proof}

\begin{remark}
The argument treats quotes as exact. It is not a statement about observation error, and it is not a statement about the far tails outside \([A,B]\). The equivalent kernel of \eqref{eq:twice} takes negative values, so the twiced estimator is not necessarily a density.
\end{remark}

\subsection{Numerical Illustration}
\label{sec:heston}

The data-generating process is the stochastic-volatility model of \citet{heston1993} with the equity parameters of \citet{bakshi1997}: \(S_0=100\), \(r=0.05\), \(q=0\), \(T=1/4\), \(\kappa=1.15\), \(\theta=0.04\), \(\sigma=0.39\), \(\rho=-0.64\), \(v_0=0.04\). The c.f.\ of \(\log S_T\) uses the little-trap branch of \citet{albrecher2007}. Call prices are the damped integrand of \citet{carrmadan1999} with damping \(\alpha=1.5\), evaluated by the trapezoid rule on \(4096\) frequencies in \((0,250]\); puts follow by parity. There is no Monte Carlo noise. The true density is Fourier inversion of the same c.f.\ on the same frequency grid,
\begin{equation}
\label{eq:true-q}
f_{\Q}(K)
\;=\;
\frac1{K\pi}\int_0^\infty
\mathrm{Re}\bigl[e^{-\mathrm{i}u\log K}\varphi(u)\bigr]\,\mathrm{d}u,
\end{equation}
where \(\varphi(u)=\E^{\Q}[e^{\mathrm{i}u\log S_T}]\). Two chains are used: a dense grid of \(m=256\) equally spaced strikes on \([30,220]\) (spacing \(\delta=0.745\)) and a sparse grid of \(m=64\) equally spaced strikes on \([70,140]\) (spacing \(\delta=1.11\)). Integrated squared error (ISE) is \(\int_{60}^{150}(\hat f_{\Q}-f_{\Q})^2\,\mathrm{d}K\), evaluated by the trapezoid rule on \(401\) equally spaced points. Table~\ref{tab:heston-tail} reports ISE for the four nested estimators at the oracle \(h^{\ast}\), at the mesh rule \(h=1.2\delta\) of \eqref{eq:h-mesh}, and at the density-scale rule \eqref{eq:h-den}. The oracle minimizes the integral on a geometric grid in \(h\).

\begin{table}[H]
\centering
\caption{Integrated squared error \(\int_{60}^{150}(\hat f_{\Q}-f_{\Q})^2\,\mathrm{d}K\) on exact Heston quotes. Mesh \(h\) is \eqref{eq:h-mesh}. The last pair of columns is \eqref{eq:h-den}. Baseline is \eqref{eq:fhat-Q} on the quoted strikes. Tail completion uses the wings of Section~\ref{sec:tails}. Tail+Midpoint is \eqref{eq:headline}. Tail+Mid+Twice is \eqref{eq:twice}.}
\label{tab:heston-tail}
\resizebox{\textwidth}{!}{%
\begin{tabular}{ll
  S[table-format=2.2]
  S[table-format=1.2e-1, scientific-notation=true]
  S[table-format=2.2]
  S[table-format=1.2e-1, scientific-notation=true]
  S[table-format=2.2]
  S[table-format=1.2e-1, scientific-notation=true]}
\toprule
\multicolumn{1}{c}{Chain} & \multicolumn{1}{c}{Estimator} &
  {\(h^{\ast}\)} & {Oracle ISE} & {Mesh \(h\)} & {Mesh ISE} & {Den.\ \(h\)} & {Den.\ ISE} \\
\midrule
Dense & Baseline \eqref{eq:fhat-Q}
  & 0.63 & 2.67e-5 & 0.89 & 2.71e-5 & 3.41 & 1.45e-4 \\
Dense & Tail completion
  & 0.63 & 2.67e-5 & 0.89 & 2.71e-5 & 3.41 & 1.45e-4 \\
Dense & Tail+Midpoint \eqref{eq:headline}
  & 0.63 & 2.27e-7 & 0.89 & 8.23e-7 & 3.41 & 1.22e-4 \\
Dense & Tail+Mid+Twice \eqref{eq:twice}
  & 0.72 & 3.44e-9 & 0.89 & 5.00e-9 & 3.41 & 8.24e-6 \\
\midrule
Sparse & Baseline \eqref{eq:fhat-Q}
  & 26.83 & 1.07e-2 & 1.33 & 2.87e2 & 4.51 & 7.42e0 \\
Sparse & Tail completion
  & 1.10 & 6.24e-5 & 1.33 & 6.34e-5 & 4.29 & 3.16e-4 \\
Sparse & Tail+Midpoint \eqref{eq:headline}
  & 0.96 & 4.39e-6 & 1.33 & 5.95e-6 & 4.29 & 2.68e-4 \\
Sparse & Tail+Mid+Twice \eqref{eq:twice}
  & 2.14 & 2.07e-6 & 1.33 & 3.24e-6 & 4.29 & 3.09e-5 \\
\bottomrule
\end{tabular}%
}
\end{table}

\begin{figure}[H]
\centering
\includegraphics[width=0.98\textwidth]{figures/ccdf_heston_rnd.pdf}
\caption{Heston density of \(S_T\) on \([60,150]\). Solid: Tail+Midpoint \eqref{eq:headline} at mesh \(h=1.2\delta\). Dashed: Tail+Mid+Twice \eqref{eq:twice}. Sparse quoted strikes are marked on the true density. Curves are truncated at zero.}
\label{fig:heston-rnd}
\end{figure}

On the dense chain \(K_1=30\) lies well to the left of the ISE window and \(C_m=0\); completing the tails does not change the baseline digits, and the mesh rule already matches that oracle (ratio \(1.02\)). Midpoints cut mesh ISE by a factor of thirty, to \(8.23\times 10^{-7}\) against an oracle \(2.27\times 10^{-7}\) (ratio \(3.6\)). Twicing then reaches mesh ISE \(5.00\times 10^{-9}\) against an oracle \(3.44\times 10^{-9}\) (ratio \(1.45\)), and is visually indistinguishable from the true density in Figure~\ref{fig:heston-rnd}. On the sparse chain \(K_1=70\) lies inside the window and \(P(K_1)=0.69\). Placing that mass at \(K_1\) makes the mesh rule unusable (ISE \(287\)). The left line \eqref{eq:left} with spacing equal to the quoted mesh brings c\`adl\`ag mesh ISE to \(6.34\times 10^{-5}\). Midpoints on the completed support are Tail+Midpoint: mesh ISE \(5.95\times 10^{-6}\) against oracle \(4.39\times 10^{-6}\) (ratio \(1.36\)). Twicing at the mesh \(h\) then reaches \(3.24\times 10^{-6}\) against an oracle \(2.07\times 10^{-6}\) (ratio \(1.56\)).

The density-scale rule \eqref{eq:h-den} is larger than the mesh on both chains (\(h=3.41\) on the dense grid, \(h=4.29\) on the sparse grid after the right wing). On exact dense quotes it oversmooths: twiced ISE \(8.24\times 10^{-6}\) against mesh \(5.00\times 10^{-9}\). On the sparse chain it is larger than both the twiced oracle and twicing at the mesh (\(3.09\times 10^{-5}\) against \(2.07\times 10^{-6}\) and \(3.24\times 10^{-6}\)). It is the rule used on listed chains in Section~\ref{sec:emp}, where the quotes are noisy.

Table~\ref{tab:vg-tail} repeats the design for the variance-gamma process of \citet{madan1998}. The process parameters are those of \citet{carrmadan1999}: \(\sigma=0.12\), \(\nu=0.2\), \(\theta=-0.14\). Spot, rate, dividend yield, and maturity are as above, as are the two strike grids and the ISE window. Variance gamma is a pure-jump L\'evy process. The true density of \(\log S_T\) is a modified Bessel function \citep{madan1998}; the density of \(S_T\) is that formula times \(1/K\). Calls are the discounted payoff of that density, evaluated by the trapezoid rule. There is no Monte Carlo noise.

\begin{table}[H]
\centering
\caption{Integrated squared error \(\int_{60}^{150}(\hat f_{\Q}-f_{\Q})^2\,\mathrm{d}K\) on exact variance-gamma quotes. Bandwidth columns as in Table~\ref{tab:heston-tail}.}
\label{tab:vg-tail}
\resizebox{\textwidth}{!}{%
\begin{tabular}{ll
  S[table-format=2.2]
  S[table-format=1.2e-1, scientific-notation=true]
  S[table-format=2.2]
  S[table-format=1.2e-1, scientific-notation=true]
  S[table-format=2.2]
  S[table-format=1.2e-1, scientific-notation=true]}
\toprule
\multicolumn{1}{c}{Chain} & \multicolumn{1}{c}{Estimator} &
  {\(h^{\ast}\)} & {Oracle ISE} & {Mesh \(h\)} & {Mesh ISE} & {Den.\ \(h\)} & {Den.\ ISE} \\
\midrule
Dense & Baseline \eqref{eq:fhat-Q}
  & 0.63 & 2.14e-4 & 0.89 & 2.42e-4 & 2.18 & 8.57e-4 \\
Dense & Tail completion
  & 0.63 & 2.14e-4 & 0.89 & 2.42e-4 & 2.18 & 8.57e-4 \\
Dense & Tail+Midpoint \eqref{eq:headline}
  & 0.54 & 1.69e-5 & 0.89 & 5.55e-5 & 2.18 & 7.18e-4 \\
Dense & Tail+Mid+Twice \eqref{eq:twice}
  & 0.63 & 2.88e-6 & 0.89 & 9.96e-6 & 2.18 & 2.18e-4 \\
\midrule
Sparse & Baseline \eqref{eq:fhat-Q}
  & 26.83 & 2.80e-2 & 1.33 & 2.87e2 & 2.88 & 2.86e1 \\
Sparse & Tail completion
  & 0.84 & 4.67e-4 & 1.33 & 5.66e-4 & 2.70 & 1.53e-3 \\
Sparse & Tail+Midpoint \eqref{eq:headline}
  & 0.84 & 5.23e-5 & 1.33 & 1.91e-4 & 2.70 & 1.26e-3 \\
Sparse & Tail+Mid+Twice \eqref{eq:twice}
  & 0.84 & 1.17e-5 & 1.33 & 4.45e-5 & 2.70 & 4.44e-4 \\
\bottomrule
\end{tabular}%
}
\end{table}

\begin{figure}[H]
\centering
\includegraphics[width=0.98\textwidth]{figures/ccdf_vg_rnd.pdf}
\caption{Variance-gamma density of \(S_T\) on \([60,150]\). Solid: Tail+Midpoint \eqref{eq:headline} at mesh \(h=1.2\delta\). Dashed: Tail+Mid+Twice \eqref{eq:twice}. Sparse quoted strikes are marked on the true density. Curves are truncated at zero.}
\label{fig:vg-rnd}
\end{figure}

On the dense chain the ranking is the same as under Heston. Completing the tails does not change the baseline digits. Midpoints cut mesh ISE from \(2.42\times 10^{-4}\) to \(5.55\times 10^{-5}\). Twicing then reaches \(9.96\times 10^{-6}\) against an oracle \(2.88\times 10^{-6}\) (ratio \(3.5\)), and tracks the peak in Figure~\ref{fig:vg-rnd}. The density of \(S_T\) is more peaked than under Heston, and the absolute errors are larger. On the sparse chain the c\`adl\`ag mesh rule is again unusable (ISE \(287\)). Tail+Midpoint mesh ISE is \(1.91\times 10^{-4}\); twicing at that \(h\) lowers it to \(4.45\times 10^{-5}\) against an oracle \(1.17\times 10^{-5}\). The density-scale rule oversmooths on both chains.

\section{Empirical Applications}
\label{sec:emp}

\subsection{Data and Methodology}
\label{sec:data}

The listed comparison uses European, AM-settled index options from Cboe's delayed-quotes feed (\texttt{cdn.cboe.com}). The S\&P~500 (SPX) snapshot is 6~September 2026, 17:04 Eastern (28{,}024 listed contracts across 55 expiries); the Nasdaq-100 (NDX) snapshot is 8~September 2026, 17:17 Eastern (14{,}650 contracts across 47 expiries); the Russell~2000 (RUT) snapshot is 19~September 2026, 03:44 Eastern (11{,}588 listed contracts). Only the standard monthly roots SPX, NDX, and RUT are kept, not weekly SPXW, NDXW, or RUTW. Time to expiry is \(T=d/365.25\) with \(d\) the calendar day count from the snapshot date. The discount rate is \(r=0.04\).

A quote is retained if the bid is positive, the ask is at least the bid, the mid is at least \(0.25\), and the spread is at most \(\max(15,0.5\times\mathrm{mid})\). Among remaining put--call pairs, the forward is the median of \(K+e^{rT}(C-P)\) over strikes with \(|K-K_{\mathrm{atm}}|\le 50\) on SPX, where \(K_{\mathrm{atm}}\) is the pair closest to the cash index; on NDX and RUT the band is \(0.65\%\) of the spot. Dividend yield is \(q=r-\log(F/S_0)/T\). The working slice is out of the money (OTM): the put mid if \(K\le F\), the call mid if \(K\ge F\), the missing side filled by parity. Puts with \(K<0.7F\) and mid below \(0.50\) are then dropped; out-of-the-money calls with mid at least \(0.25\) are kept. Table~\ref{tab:data} records the filter counts. In \eqref{eq:headline} and \eqref{eq:twice}, \(\tilde C=C/(S_0e^{-qT})\) and \(F=S_0e^{(r-q)T}\).

The December 2026 S\&P~500 AM quarterly (\(T=0.282\), \(S_0=7718.60\), \(F=7786\), \(q=0.92\%\)) appears in Figure~\ref{fig:listed}. After the screens it has \(n=393\) strikes on \([0.31F,1.36F]\), of which \(77\%\) are puts. Near the money the strike mesh is \(5\) (\(0.06\%\) of \(F\)); it is \(10\)--\(25\) on the put wing and up to \(200\) beyond \(1.2F\). Median out-of-the-money mids are \(51\) for puts and \(77\) for calls. Implied volatility \(\hat\sigma_{\mathrm{ATM}}=0.141\) is the median Black--Scholes implied volatility of contracts with \(|K-F|\le 0.03F\). The linear call tail of Section~\ref{sec:tails} is appended (\(C_m=0.25\) by the mid floor of the filter, \(n_R=80\) from \eqref{eq:nR}), which produces \(n=473\) masses for \eqref{eq:h-den} and \(h=180\). The March 2027 AM quarterly (\(T=0.531\), \(F=7864\), \(q=0.49\%\), \(\hat\sigma_{\mathrm{ATM}}=0.147\)) filters to \(n=222\) strikes on \([0.18F,1.50F]\), with a \(25\)-point mesh at the money; \eqref{eq:nR} gives \(n_R=52\) extra points and \(h=291\).

The Nasdaq-100 December 2026 AM quarterly (\(T=0.277\), \(S_0=29596\), \(F=29898\), \(q=0.33\%\), \(\hat\sigma_{\mathrm{ATM}}=0.204\)) is built with the same screens. It filters to \(n=252\) strikes on \([0.23F,1.51F]\), \(73\%\) puts; the at-the-money mesh is \(25\), and gaps of \(1000\) appear beyond \(1.15F\). Density-scale \(h=1082\) on \(n=308\) masses after the call wing (\eqref{eq:nR} gives \(n_R=56\)).

The Russell~2000 December 2026 AM quarterly (\(T=0.246\), \(S_0=2860\), \(F=2883\), \(q=0.84\%\), \(\hat\sigma_{\mathrm{ATM}}=0.185\)) filters to \(n=97\) strikes on \([0.43F,1.28F]\), \(58\%\) puts; the at-the-money mesh is \(10\). Density-scale \(h=97.8\) on \(n=193\) masses after the call wing (\eqref{eq:nR} gives \(n_R=96\)).

Out-of-the-money RMSE uses the quoted put for \(K\le F\) and the quoted call for \(K\ge F\). Hold-out is the average of two even/odd splits (train on even strikes and score odd, then the reverse), with \(h\) recomputed from the training \(n\). We compare the method with four competitors.

\begin{table}[H]
\centering
\caption{Construction of the working slices. ``At expiry'' counts listed puts and calls on that root and maturity with a positive bid and ask at least the bid. ``Screens'' applies the mid and spread filters. ``Working OTM'' is the slice used by the estimators.}
\label{tab:data}
{\small
\begin{tabular}{@{}l *{5}{S[table-format=3.0]}@{}}
\toprule
\multicolumn{1}{@{}c}{Slice} & {At expiry} & {Screens} & {Unique \(K\)} & {Working OTM} & {Puts} \\
\midrule
SPX 18 Dec 2026 & 819 & 812 & 416 & 393 & 303 \\
SPX 19 Mar 2027 & 453 & 452 & 228 & 222 & 164 \\
NDX 18 Dec 2026 & 555 & 531 & 279 & 252 & 185 \\
RUT 18 Dec 2026 & 217 & 216 & 115 & 97 & 56 \\
\bottomrule
\end{tabular}
}
\end{table}

\subsection{Estimators}
\label{sec:emp-est}

\paragraph{Ours}
Start from quoted strikes \(0<K_1<\cdots<K_m\) and calls \(C_1,\ldots,C_m\). Set \(\tilde C_i=C_i/(S_0e^{-qT})\), clip to \([0,1]\), and replace \(\tilde C\) by its running minimum so that \(P_i=1-\tilde C_i\) is a cdf, \(0\le P_1\le\cdots\le P_m\le 1\). Then:

\textit{Right wing.} If \(C_m>0\), let
\[
s
\;=\;
\mathrm{clip}\Bigl(\frac{C_m-C_{m-1}}{K_m-K_{m-1}},\,-e^{-rT},\,0\Bigr),
\]
and \(K_{\mathrm{end}}=\min\bigl(\max(K_m+C_m/(-s),\,1.01 K_m),\,8F\bigr)\). Append \(n_R\) equally spaced strikes in \((K_m,K_{\mathrm{end}}]\) from \eqref{eq:nR}, with
\[
C(K)=\max\bigl(C_m+s(K-K_m),0\bigr),
\qquad
P(K)=1-C(K)/(S_0e^{-qT}).
\]
If \(C_m=0\), skip this wing.

\textit{Left wing.} Remove the jump \(P_1\) at \(K_1\). Let \(\delta\) be the median quoted gap and \(n_L=\mathrm{round}(K_1/\delta)\). Put \(n_L\) equal jumps of size \(P_1/n_L\) at
\[
K_j^{L}
\;=\;
j\,K_1/n_L,
\qquad
j=1,\ldots,n_L.
\]
That is the discretization of \eqref{eq:left}.

\textit{Midpoints.} Let \(0=K_0^{\ast}<K_1^{\ast}<\cdots<K_N^{\ast}\) be the resulting ordered knots (left-wing points, quoted strikes from \(K_2\) onwards, right-wing points) and \(\Delta_j\) the jump on \((K_{j-1}^{\ast},K_j^{\ast}]\). Set \(\tilde K_j=\tfrac12(K_{j-1}^{\ast}+K_j^{\ast})\). Form \eqref{eq:headline} and \eqref{eq:P-hat-mid}, then \eqref{eq:twice}. Clip \(\hat P^{\mathrm{tw}}\) to \([0,1]\). Calls are \(\hat C=S_0e^{-qT}(1-\hat P^{\mathrm{tw}})\), puts by parity. The listed rule is density-scale \eqref{eq:h-den}. 

\paragraph{\citeauthor{yatchew2006}}
The quoted calls \(C=(C_1,\ldots,C_m)\) are replaced by the nearest vector \(m=(m_1,\ldots,m_m)\) that is decreasing and convex with consecutive slopes in \([-e^{-rT},0]\), and that lies above intrinsic value \(e^{-rT}\max(F-K_i,0)\). The implementation parameterizes that cone as follows. Let \(\Delta K_i=K_{i+1}-K_i\). Write \(m_1\) free and
\[
m_{i+1}
\;=\;
m_i+s_i\,\Delta K_i,
\qquad
s_i
\;=\;
-e^{-rT}+\sum_{j=1}^{i}v_j,
\]
with \(v_j\ge 0\) and \(\sum_j v_j\le e^{-rT}\). Then the slopes increase from \(-e^{-rT}\) to at most \(0\), so \(m\) is decreasing and convex with legal digital bounds. The objective is \(\tfrac12\|m-C\|^2\); there is no roughness penalty (\(\lambda=0\) in their notation). The resulting nonlinear program is solved by sequential least-squares programming (SLSQP), started from a pool-adjacent-violators projection of the clipped discrete slopes of \(C\). After the fit, \(m\) is clipped above intrinsic and the slopes are re-isotonized. If \(C\) already lies in the cone, the projection is the identity. Prices in Table~\ref{tab:listed} are this list at the quoted strikes, with puts by parity. Mid quotes on the S\&P~500 and Nasdaq-100 chains violate butterfly arbitrage, so the projection moves them and the in-sample RMSE is not zero. The mass column uses \(e^{rT}\) times the second difference of \(m\).

\paragraph{\citeauthor{aitsahalia2003}}
\citet{aitsahalia2003} first replace the quoted calls by a decreasing convex list, then apply a local polynomial to the constrained values. The constrained list is the same cone as \citeauthor{yatchew2006} above. The call interpolant is a local linear of that list, with Gaussian weights and bandwidth \(1.06\,F\,\hat\sigma\sqrt{T}\,n^{-1/5}\). The density is \(e^{rT}\) times twice the quadratic coefficient of a local cubic of the same constrained list, with bandwidth \(0.9\,F\,\hat\sigma\sqrt{T}\,n^{-1/9}\), the usual rate for a second derivative. \(\hat\sigma\) is the median implied volatility of the constrained calls. Prices in Table~\ref{tab:listed} are the local-linear interpolant at the quoted strikes, with puts by parity.

\paragraph{\citeauthor{bondarenko2003} PCA}
The density is required to be a mixture of lognormals that all share a common volatility \(h\), with positive weights adding to one. Let \(\hat\sigma\) be the median Black--Scholes implied volatility of the slice and \(s=\hat\sigma\sqrt{T}\). Centers \(z_1,\ldots,z_{41}\) are equally spaced on \([\log F-4s,\log F+4s]\), and \(h=0.65s\). The density is
\begin{equation}
\label{eq:pca-q}
\hat q_{\mathrm{PCA}}(K)
\;=\;
\sum_{j=1}^{41} a_j\,\frac{1}{Kh\sqrt{2\pi}}
\exp\Bigl(-\frac{(\log K-z_j)^2}{2h^2}\Bigr),
\qquad
a_j\ge 0,\quad
\sum_{j}a_j=1.
\end{equation}
Each component prices a call in closed form (Black--Scholes with forward \(e^{z_j+h^2/2}\)). The weights are fitted by nonnegative least squares to the observed calls, with two additional penalties that keep the weights summing to one and the mean equal to the forward. The estimator is a density by construction. With \(41\) centers and \(h\) of order \(s\), the approximation does not become finer as the strike list becomes denser: it is a global mixture, not a local smoother. Prices in Table~\ref{tab:listed} are the integral of \(\max(\hat q_{\mathrm{PCA}},0)\) against the put and call payoffs on \([\max(50,0.2F),2.4F]\).

\paragraph{\citeauthor{priestley1972}}
The listed implementation is the cubic interpolant of Section~\ref{sec:naive}: a clamped cubic spline through the quotes, extended by linear wings of slope in \([-e^{-rT},0]\), then convolved with a Gaussian kernel. The density is \eqref{eq:pc-cubic}. Prices in Table~\ref{tab:listed} are \((C^{\mathrm{cs}}\ast\kappa_h)\) at the quoted strikes, with puts by parity, not the unsmoothed spline. The bandwidth is \(1.06\,F\,\hat\sigma_{\mathrm{ATM}}\sqrt{T}\,n^{-1/9}\), the usual rate for a second derivative; \(n\) is the number of quoted strikes. Hold-out refits the spline on the training knots.

\begin{table}[H]
\centering
\caption{Out-of-the-money repricing. Ours is \eqref{eq:twice} at the density-scale bandwidth \eqref{eq:h-den} (\(h=180\) on SPX December, \(h=291\) on SPX March, \(h=1082\) on NDX December, \(h=97.8\) on RUT December). \citeauthor{yatchew2006} \(\lambda=0\) from the constrained call; \citeauthor{aitsahalia2003} is priced from its call interpolant; PCA from its density. \citeauthor{priestley1972} is the cubic interpolant \eqref{eq:pc-cubic} at \(1.06\,F\,\hat\sigma_{\mathrm{ATM}}\sqrt{T}\,n^{-1/9}\) (\(h=318\) on SPX December, \(h=491\) on SPX March, \(h=1841\) on NDX December, \(h=169\) on RUT December), priced from the kernel-smoothed cubic. OTM is put RMSE for \(K\le F\) and call RMSE for \(K\ge F\). Mass is \(\int\max(\hat q,0)\) on \([\max(50,0.2F),2.4F]\). Hold-out is the average OTM RMSE of the two even/odd splits.}
\label{tab:listed}
\resizebox{\textwidth}{!}{%
\begin{tabular}{ll *{5}{S[table-format=3.2]}}
\toprule
\multicolumn{1}{c}{Chain} & \multicolumn{1}{c}{Method} &
  {RMSE OTM} & {Puts} & {Calls} & {Hold-out} & {Mass} \\
\midrule
SPX Dec & Ours
  & 0.49 & 0.31 & 0.85 & 0.95 & 1.01 \\
SPX Dec & \citeauthor{yatchew2006}
  & 2.25 & 2.09 & 2.72 & 4.50 & 1.00 \\
SPX Dec & \citeauthor{aitsahalia2003}
  & 7.95 & 6.28 & 11.98 & 11.29 & 1.02 \\
SPX Dec & PCA
  & 10.99 & 10.10 & 13.56 & 10.62 & 1.00 \\
SPX Dec & \citeauthor{priestley1972}
  & 20.48 & 16.27 & 30.66 & 23.66 & 1.00 \\
\midrule
SPX Mar & Ours
  & 0.95 & 0.64 & 1.53 & 1.90 & 1.02 \\
SPX Mar & \citeauthor{yatchew2006}
  & 6.91 & 6.23 & 8.54 & 9.82 & 1.00 \\
SPX Mar & \citeauthor{aitsahalia2003}
  & 17.63 & 14.19 & 24.91 & 28.37 & 1.02 \\
SPX Mar & PCA
  & 31.33 & 27.43 & 40.36 & 30.20 & 1.00 \\
SPX Mar & \citeauthor{priestley1972}
  & 28.63 & 19.85 & 44.98 & 33.12 & 1.00 \\
\midrule
NDX Dec & Ours
  & 6.45 & 7.19 & 3.68 & 8.74 & 1.15 \\
NDX Dec & \citeauthor{yatchew2006}
  & 20.44 & 13.24 & 32.99 & 53.81 & 1.00 \\
NDX Dec & \citeauthor{aitsahalia2003}
  & 54.24 & 42.87 & 77.41 & 23.71 & 1.03 \\
NDX Dec & PCA
  & 25.14 & 29.26 & 3.46 & 22.85 & 1.00 \\
NDX Dec & \citeauthor{priestley1972}
  & 135.69 & 115.99 & 179.18 & 156.52 & 1.00 \\
\midrule
RUT Dec & Ours
  & 0.40 & 0.34 & 0.48 & 0.96 & 1.01 \\
RUT Dec & \citeauthor{yatchew2006}
  & 0.01 & 0.01 & 0.00 & 5.19 & 1.00 \\
RUT Dec & \citeauthor{aitsahalia2003}
  & 5.58 & 4.60 & 6.69 & 4.40 & 0.99 \\
RUT Dec & PCA
  & 1.96 & 2.58 & 0.07 & 1.73 & 1.00 \\
RUT Dec & \citeauthor{priestley1972}
  & 13.79 & 11.43 & 16.47 & 15.92 & 1.00 \\
\bottomrule
\end{tabular}%
}
\end{table}

The quotes are not in the \citeauthor{yatchew2006} cone on SPX and NDX (in-sample RMSE \(2.25\) on SPX December, \(6.91\) on SPX March, \(20.44\) on NDX December). That residual is a floor for any decreasing convex interpolant of \(C\): second differences of the projection are a discrete nonnegative measure with mass \(1.00\), but not a curve. On RUT December the mids already lie in the cone (projection residual \(0.01\)). With \(\lambda=0\) the projection is then the identity in sample, and hold-out linearly interpolates the training knots: in-sample RMSE \(0.01\), hold-out \(5.19\). The kernel smoother has in-sample RMSE \(0.40\) and hold-out \(0.96\). PCA is a density by construction and is fitted to prices; with \(41\) shared-width lognormals it cannot track a few hundred noisy strikes, so out-of-the-money RMSE is \(10.99\) on SPX December. \citeauthor{aitsahalia2003} stay inside the cone and then smooth, so they cannot beat the floor either. Prices in Table~\ref{tab:listed} for the estimator of this paper are the interpolant \(\hat C=S_0e^{-qT}(1-\hat P^{\mathrm{tw}})\). That interpolant is not required to be convex, and its out-of-the-money RMSE lies below the floor (\(0.49\) on SPX December, \(0.95\) on SPX March, \(6.45\) on NDX December), in sample and on hold-out. The same kernel yields the smooth estimate of \(f_{\Q}\) in the left panel of Figure~\ref{fig:listed}. The equivalent kernel of \eqref{eq:twice} takes negative values; clipped mass is \(1.01\)--\(1.02\) on S\&P~500 and Russell~2000 and \(1.15\) on Nasdaq-100. PCA is the more symmetric density; the Gaussian estimator puts more mass below \(F\). The \citeauthor{yatchew2006} interpolant appears on the implied-volatility panel only. The \citeauthor{priestley1972} estimator kernels the cubic interpolant of the call levels at the second-derivative bandwidth. That \(h\) is larger than \eqref{eq:h-den}, and the prices are a kernel average of a curve that is large deep in the money, so out-of-the-money RMSE is \(20.48\) on SPX December and \(135.69\) on NDX December, in sample and on hold-out.

\begin{figure}[H]
\centering
\includegraphics[width=0.98\textwidth]{figures/ccdf_listed.pdf}
\caption{SPX, NDX, and RUT, 18 December 2026 AM-settled chains. Left: \(\max(\hat f_{\Q},0)\) from Ours \eqref{eq:twice} at density-scale \(h\), \citeauthor{aitsahalia2003}, PCA, and \citeauthor{priestley1972} \eqref{eq:pc-cubic}. Right: out-of-the-money implied volatility from Ours, \citeauthor{yatchew2006}, PCA, and \citeauthor{priestley1972}. The dashed line is the forward.}
\label{fig:listed}
\end{figure}

\section{Conclusion}
\label{sec:conc}

The second strike derivative of a European call is the risk-neutral density of the underlying at expiry. On a finite strike list that derivative is not an estimator. This paper rescales the calls to a complementary distribution function, convolves the empirical jumps with a Gaussian kernel, and inverts the associated characteristic function in closed form. Completing the unquoted tails, placing masses at cell centers, and combining two kernel widths are quadrature corrections. The resulting estimator is in-fill consistent for exact quotes, and on a dense exact Heston chain the mesh rule already attains an integrated squared error of \(5.0\times 10^{-9}\). The same ranking of the four nested estimators holds on a variance-gamma chain with the same strike grids. On listed S\&P~500, Nasdaq-100, and Russell~2000 strips the same kernel supplies a smooth estimate of \(f_{\Q}\) and a call interpolant. At a density-scale bandwidth that interpolant reprices out-of-the-money mids on the S\&P~500 and Nasdaq-100 chains more closely than \citeauthor{bondarenko2003}'s mixture, than the no-arbitrage projection of \citeauthor{yatchew2006}, than \citeauthor{aitsahalia2003}, and than a \citeauthor{priestley1972} kernel of the call, including on an even/odd hold-out averaged over both directions.

\section*{Code and Data Availability}
A Python implementation of \eqref{eq:twice} and a replication script are available at \url{https://github.com/s-broda/RND}.

\bibliographystyle{elsarticle-harv}
\bibliography{refs}
\end{document}
